% Einheit 1 von 2 – Jahreszins, Monats- und Tageszins (bank/zinsrechnung/e1.jsonl, zusammenbau v0.1)
%% TODO Zweigzeile Teil 1: was hier gelernt wird (Satz in Schülersprache aus dem Einheitstitel) – Einheit 1
\einheitenkopf[e1]{Einheit 1 von 2 · Jahreszins, Monats- und Tageszins}
\zweigzeile{ab Kl. 8 (am Gymnasium ab Kl. 7) · P10}
\setcounter{aufgabe}{10}

%% TODO \verfahren{…}: Verfahrensüberschrift in Sachsprache fehlt in der Bank (2.3 b)
%% TODO Titel: Ich-kann-Satz fehlt in der Bank (Platzhalter: Kettenname) – Nr. 11
%% TODO Anweisung: ein Satz über der ersten Teilaufgabe fehlt in der Bank – Nr. 11
\begin{aufgabe}{Zinsen berechnen}
%% TODO \rechenplatz{n}: Schrittzahl des Grundfalls fehlt in der Bank (nur bei fester Schreibform, 2.3 b)
\begin{teile}
\teil Lena legt 850\,€ für ein Jahr auf ein Sparbuch, die Bank zahlt 2\,\% Zinsen. Wie viel Euro Zinsen bekommt sie? Was ist gesucht? Kreuze an, rechne nicht. \\ \kreuz{das Kapital (Grundwert)} \\ \kreuz{der Zinssatz (Prozentsatz)} \\ \kreuz{die Zinsen (Prozentwert)}
\teil Sparbuch: 300\,€ zu 2\,\% – wie viel Euro Zinsen nach einem Jahr? \leerfeld[€]
\teil Rechne über 1\,\%: 600\,€ zu 4\,\% – wie viel Euro Zinsen im Jahr? \leerfeld[€]
\teil Jana legt 700\,€ für ein Jahr zu 2\,\% an. Wie viel Euro hat Jana nach einem Jahr auf dem Konto? \leerfeld[€]
\teil 900\,€ Kapital bringen in einem Jahr 27\,€ Zinsen. Wie hoch ist der Zinssatz? \leerfeld[\%]
\teil Guthaben 4\,000\,€, Zinsen in einem Jahr 60\,€ – welcher Zinssatz? \leerfeld[\%]
\teil 21\,€ Zinsen im Jahr sind 3\,\% vom Kapital. Wie viel Euro Kapital? \leerfeld[€]
\teil Weise mit der Tabelle nach, dass die Bank 3\,\% Zinsen im Jahr zahlt.

\sachtabelle{lrrr}{Buchung & Einzahlung & Zinsen & Guthaben}{Jahresbeginn & 300,00\,€ & – & 300,00\,€\\ Jahresende & – & 9,00\,€ & 309,00\,€}
\teil 600\,€ zu 2\,\%, die Zinsen werden jedes Jahr ausgezahlt. Wie viel Euro Zinsen in 3 Jahren? \leerfeld[€]
\teil Für ein Motorroller leiht sich Familie Roth 1\,800\,€ für ein Jahr zu 6\,\%. Wie viel Euro muss sie nach einem Jahr zurückzahlen? \leerfeld[€]
\teil 3\,000\,€ für ein Jahr: Bank A zahlt 2\,\%, Bank B 1,5\,\% und 20\,€ Startbonus. Welche Bank bringt mehr? Bank \leerfeld
\teil Ein Sparbetrag von 2\,600\,€ wird ein Jahr lang mit 2\,\% pro Jahr verzinst. Wie viel Euro Zinsen gibt es nach einem Jahr? \hfill (P10 2015 OS) \\ \leerfeld[€]
\end{teile}
\end{aufgabe}

%% TODO \verfahren{…}: Verfahrensüberschrift in Sachsprache fehlt in der Bank (2.3 b)
%% TODO Titel: Ich-kann-Satz fehlt in der Bank (Platzhalter: Kettenname) – Nr. 12
%% TODO Anweisung: ein Satz über der ersten Teilaufgabe fehlt in der Bank – Nr. 12
\begin{aufgabe}{Monats- und Tageszinsen}
%% TODO \rechenplatz{n}: Schrittzahl des Grundfalls fehlt in der Bank (nur bei fester Schreibform, 2.3 b)
\begin{teile}
\teil Geld liegt von Anfang März bis Ende August auf dem Konto. Welche Laufzeit? Kreuze an, rechne nicht. \\ \kreuz{ein Jahr – nicht teilen} \\ \kreuz{Monate – durch zwölf teilen} \\ \kreuz{Tage – durch dreihundertsechzig teilen}
\teil 2\,400\,€ zu 2\,\% für 5 Monate – wie viel Euro Zinsen? \leerfeld[€]
\teil 7\,200\,€ zu 2\,\% für 40 Tage, Bankjahr – wie viel Euro Zinsen? \leerfeld[€]
\teil 5\,400\,€ zu 4\,\% vom 11. März bis zum 26. Mai, jeder Monat mit dreißig Tagen – wie viel Euro Zinsen? \leerfeld[€]
\teil 3\,000\,€ bringen in 4 Monaten 20\,€ Zinsen. Welcher Zinssatz? \leerfeld[\%]
\teil Nora hat 4\,200\,€ aus ihrem Nebenjob gespart. Sie legt das Geld von Anfang Mai bis Ende Oktober auf ein Tagesgeldkonto mit 2,5\,\% Zinsen im Jahr. Wie viel Euro Zinsen bekommt sie? Antwort: \leerfeld
\end{teile}
\end{aufgabe}

%% TODO Titel: Ich-kann-Satz fehlt in der Bank (Platzhalter: Kettenname) – Nr. 13
%% TODO Anweisung: ein Satz über der ersten Teilaufgabe fehlt in der Bank – Nr. 13
\begin{aufgabe}{Überschlag: 1 \% von … , dann grob mal p}
\begin{teile}
\teil Überschlage: 3\,980\,€ zu 2\,\% – etwa wie viel Euro Zinsen im Jahr? ≈ \leerfeld[€]
\end{teile}
\end{aufgabe}

%% TODO Titel: Ich-kann-Satz fehlt in der Bank (Platzhalter: Kettenname) – Nr. 14
%% TODO Anweisung: ein Satz über der ersten Teilaufgabe fehlt in der Bank – Nr. 14
\begin{aufgabe}{Zinsen berechnen – Fehler finden, Begründen, Anwendung}
\begin{teile}
\teil 950\,€ zu 4\,\% – wie viel Euro Zinsen nach einem Jahr? Aylin rechnet: \rechnung{950 \cdot 0{,}04 &= 38 \\ 950 + 38 &= 988} Antwort: 988\,€. Finde den Fehler.
\teil Warum bringt 1\,\% von 20\,000\,€ mehr Zinsen als 4\,\% von 600\,€? Begründe ohne genau zu rechnen.
\teil Frau Demir leiht sich 4\,000\,€ für ein Jahr zu 6\,\%. Sie kann jeden Monat 350\,€ zurücklegen. Reicht das Gesparte nach einem Jahr für die Rückzahlung? Antwort: \leerfeld
\end{teile}
\end{aufgabe}
